\section{Introduction}

Compartmental epidemiological models are the standard tool for studying infectious disease
dynamics. In such models, a population is partitioned into mutually exclusive states
(e.g., Susceptible, Exposed, Infectious, Recovered), and transitions between them are
governed by differential equations parameterized by rates such as transmission
coefficient $\beta$ or recovery rate $\gamma$~\cite{keeling2008modeling}.
Because this structure is widely familiar, authors often treat many components as
conventional and leave them implicit, describing only deviations from a standard template.
While generally harmless within a single research group, this practice creates a
\textit{specification--validation gap}: the model described in the text is not necessarily
the one validated in equations or code, nor the one a reader can reproduce. This gap
inhibits reproducibility, reusability, and cross-disease extensibility, as seen during
the COVID-19 pandemic when models were difficult to compare without harmonizing
structures and assumptions. More broadly, published compartmental models exhibit
substantial variation in structure and notation~\cite{keeling2008modeling}.

Manual reconstruction of such models does not scale. We use large language models (LLMs)
as \emph{extraction instruments}: they ingest full-length PDFs, identify epidemiological
constructs, and produce structured representations. Rather than allowing LLMs to
\emph{generate} models, which shifts the gap into opaque behavior, we constrain them to
recover only what each paper states (explicitly or by convention), and then measure
deviation from both paper promises and an independent reference model. This turns the
specification--validation gap into a measurable property of the paper rather than an
informal judgment.

From a model-driven engineering perspective, the pipeline operates over an explicit
epidemiological metamodel instantiated as \texttt{.compmodel} artifacts. Phase~2 performs
a paper-to-model transformation producing metamodel-conformant instances, while Phase~3
performs model validation, gap analysis, and rule-based repair, aligning the approach
with core MDE concerns of model conformance, completeness, and transformation quality.

We present a three-phase pipeline that operationalizes this idea at scale.
\textbf{Phase~1} parses reference (``gold-standard'') XML models and constructs structured
representations of compartments, parameters, flows, and sensitivities.
\textbf{Phase~2} extracts a draft model from each paper by combining structured text
parsing, LLM-assisted entity extraction, and XML synthesis. We evaluate three commercial
LLM backends (OpenAI, Gemini, Claude) using a domain-aware fuzzy matching protocol that
handles epidemiological label variation.
\textbf{Phase~3} refines the draft via three-layer gap analysis, retrieval-augmented
completion from the paper text, optional LLM inference, and rule-based structural repair.
A composite \emph{completeness score} (0--100) summarizes coverage, traceability, accuracy,
and structural integrity.

We apply this pipeline to 10 peer-reviewed papers spanning 10 infectious diseases:
cholera~\cite{Fung2014Cholera}, COVID-19~\cite{TuiteE497}, dengue~\cite{dengue},
Ebola~\cite{Abbas2024Ebola}, influenza~\cite{Corbella2017MonitoringAP},
HIV~\cite{Espitia2022}, malaria~\cite{Akowe2025Malaria}, measles~\cite{Pokharel2024Measles},
tuberculosis~\cite{tuberculosis}, and Zika~\cite{Lee2016Zika}. Models are represented in
the EpiMDE platform~\cite{epimde2024}, a domain-specific language based on
\texttt{.compmodel} files, enabling consistent comparison across extracted and reference
models. The main contributions of this paper are:

\begin{enumerate}[leftmargin=*, label=\arabic*.]
  \item A model-driven measurement framework for epidemiological model completeness that
        separates paper promises, extracted content, and reference agreement, making the
        specification--validation gap a measurable artifact property.
  \item A fuzzy matching protocol for compartmental model entities that handles
        transliteration, synonym groups, and naming variation.
  \item A model completion, validation, and structural repair pipeline using
        paper-grounded retrieval and deterministic repair rules.
  \item A reproducibility-oriented completeness score measuring coverage, reference
        agreement, traceability, parameter accuracy, and structural integrity.
  \item An empirical study across 10 diseases identifying systematic extraction gaps and
        structural error patterns.
\end{enumerate}

This paper extends our earlier work~\cite{11344468}, which demonstrated LLM-based model
generation for a single SEIR disease. Here we introduce a full validation pipeline,
a completeness metric, and a cross-disease empirical analysis.